\documentclass[11pt]{article}
\usepackage[a4paper,margin=23mm]{geometry}
\usepackage{amsmath,amssymb,xcolor,hyperref,fancyhdr}
\usepackage{fontspec}
\setmainfont{Times New Roman}
\setmonofont{Times New Roman}
\usepackage[italic]{mathastext}
\setlength{\emergencystretch}{3em}
\setlength{\parskip}{7pt}\setlength{\parindent}{0pt}
\pagestyle{fancy}\fancyhf{}\setlength{\headheight}{14pt}
\lhead{FlowMLLab}\rhead{Reading and experiment guide}\cfoot{\thepage}
\begin{document}
\noindent\textbf{Week 7.3 | Reading the evidence}\par\medskip\hrule\medskip
\section*{Compare information and compute budgets together}
The main scratch fit uses 300 downstream steps; the extended scratch comparison
uses 5000. Those are different experiments. At each label count, report training
steps, optimizer, unlabelled source data, target labels and evaluation frames
beside the error. The extended control can reverse a conclusion drawn from the
short-budget curve. A pretrained initialization is not automatically better at
equal compute.

The pooled POD baseline must see the same allowed unlabelled frames if the claim
concerns the value of pretraining. Identical errors at one and two labels can
mean validation kept the same initial checkpoint, rather than that another label
contains no information. CPU training can take tens of minutes; the fair-budget
control is a substantial part of the runtime. Use the printed fresh output path
to find this run; saved historical text is not a new execution.
\paragraph{Check your understanding.}
Identify the input information, the reference, the metric denominator and the
evaluation split in one worked notebook comparison. Predict which conclusion
would change if any one of those four changed, then check your prediction.
\end{document}
